\documentclass[11pt]{article}
\usepackage[margin=1in]{geometry}
\usepackage{amsmath,amssymb,amsthm,mathtools}
\usepackage{graphicx}
\usepackage[colorlinks=true,linkcolor=blue,citecolor=blue,urlcolor=blue]{hyperref}
\usepackage{microtype}

\newtheorem{theorem}{Theorem}
\newtheorem{lemma}{Lemma}
\newtheorem{corollary}{Corollary}
\newcommand{\R}{\mathbb R}
\newcommand{\F}{\mathcal F}

\title{The motivating cusp symbol in arXiv:2601.11934\\is already locally Wiener}
\author{Candidate partial-result packet}
\date{11 August 2026}

\begin{document}
\maketitle

\noindent\textbf{Status.}
Candidate substantial partial result; likely valid; human review recommended.
It settles the explicit symbol displayed in Remark~1.5 of the source, but it
does not remove the Wiener hypothesis for every $C^m$ symbol.

\section{Source question and claimed obstruction}

Chen and Hong study nonlinear superposition operators on quantum Besov spaces
in their 2026 preprint~\cite{ChenHong}. Their Theorem~1.4 assumes, locally,
$F\in W_0(\R)\cap W_{\lceil s\rceil}(\R)$. Remark~1.5 states that when
$\lceil s\rceil=2k+1$ and
\[
    F(x)=|x|x^{2k+1},
\]
this symbol is not locally in $W_0\cap W_{2k+1}$, and presents it as the
motivating example for the regularity gap. The source passage is reproduced
below (PDF page~6).

\begin{center}
\includegraphics[width=\textwidth]{figures/open_problem_crop.png}
\end{center}

On PDF page~23, the source defines
\[
 W_n(\R)=\{G\in C^n(\R):G\in L_\infty(\R),\ F(G^{(n)})\in L_1(\R)\}.
\]
Membership is local when every smooth compact localization has this property.
Under that definition, the displayed nonmembership assertion is false.

\section{Idea of the proof}

The key observation is one-dimensional and elementary. A compactly supported
$H^1$ function has integrable Fourier transform: insert the weight
$(1+\xi^2)^{-1/2}$ and use Cauchy--Schwarz and Plancherel. For the cusp symbol,
the allegedly troublesome top derivative is just a constant multiple of
$|x|$. After a compact cutoff it is Lipschitz, hence in $H^1$. The same is
true of the localized original function, so both $W_0$ and $W_m$ conditions
hold. This also yields a useful sufficient class substantially broader than
the single example.

\section{A local Wiener criterion}

We use any standard unitary or nonunitary normalization of the Fourier
transform; only harmless constants change.

\begin{lemma}\label{lem:H1}
If $h\in H^1(\R)$, then $\widehat h\in L_1(\R)$ and
\[
 \|\widehat h\|_1\le C\|h\|_{H^1}.
\]
\end{lemma}

\begin{proof}
By Cauchy--Schwarz and Plancherel,
\begin{align*}
 \int_{\R}|\widehat h(\xi)|\,d\xi
 &\le
 \left(\int_{\R}\frac{d\xi}{1+\xi^2}\right)^{1/2}
 \left(\int_{\R}(1+\xi^2)|\widehat h(\xi)|^2\,d\xi\right)^{1/2}\\
 &\le C\|h\|_{H^1}.
\end{align*}
\end{proof}

\begin{theorem}\label{thm:criterion}
Let $m\ge1$. Suppose that $F\in C^m(\R)$ and
$F^{(m)}\in H^1_{\mathrm{loc}}(\R)$. Then
\[
             F\in_{\mathrm{loc}} W_0(\R)\cap W_m(\R).
\]
In particular, the conclusion holds whenever $F^{(m)}$ is locally Lipschitz.
\end{theorem}

\begin{proof}
Let $\chi\in C_c^\infty(\R)$ and put $G=\chi F$. The function $G$ is bounded,
compactly supported, and belongs to $C^m$. Since $m\ge1$, it is also a
compactly supported $C^1$ function, hence $G\in H^1(\R)$. Lemma~\ref{lem:H1}
gives $\widehat G\in L_1$, which is the $W_0$ condition.

Leibniz's formula gives
\[
 G^{(m)}=\sum_{j=0}^{m}\binom mj\chi^{(m-j)}F^{(j)}.
\]
The term with $j=m$ lies in $H^1$: multiplication by a smooth compactly
supported function preserves local $H^1$. If $j<m$, then $F^{(j)}$ is locally
$C^1$, so $\chi^{(m-j)}F^{(j)}$ is compactly supported $C^1$ and therefore is
in $H^1$. Thus $G^{(m)}\in H^1(\R)$. A second application of
Lemma~\ref{lem:H1} yields $\widehat{G^{(m)}}\in L_1$, which is the $W_m$
condition. Since $\chi$ was arbitrary, the membership is local.
\end{proof}

\section{The cusp in Remark 1.5}

\begin{corollary}\label{cor:cusp}
Let $m=2k+1$ be odd and set $F(x)=|x|x^m$. Then
\[
       F\in_{\mathrm{loc}}W_0(\R)\cap W_m(\R).
\]
Consequently, whenever $\lceil s\rceil=m$, this symbol already satisfies the
symbol hypothesis of Theorem~1.4 in \cite{ChenHong}; hence that theorem gives
its stated quantum Besov estimate for all of its advertised $p,q,s$ and
bounded self-adjoint inputs.
\end{corollary}

\begin{proof}
On the positive and negative half-lines,
\[
 F(x)=\begin{cases}x^{m+1},&x>0,\\-x^{m+1},&x<0.\end{cases}
\]
Differentiating $m$ times gives, on both half-lines,
\[
                F^{(m)}(x)=(m+1)!|x|.
\]
Both sides are continuous at zero, so this is the classical $m$-th derivative
on all of $\R$. Its weak derivative is $(m+1)!\operatorname{sgn}(x)$ almost
everywhere, which belongs to $L_2$ on every compact interval. Therefore
$F^{(m)}\in H^1_{\mathrm{loc}}$, and Theorem~\ref{thm:criterion} applies.
Finally $F(0)=0$, so the remaining symbol hypothesis of the source theorem is
also satisfied.
\end{proof}

\section{Verification and limitations}

The proof was checked directly against the current arXiv v2 definition of
$W_n$. It uses only Leibniz's formula, Cauchy--Schwarz, and Plancherel. The
source crop contains the whole of Theorem~1.4 and Remark~1.5, with no omitted
continuation. No computational experiment is being presented as evidence.

The result does not show that every $C^m$ function is locally in $W_m$.
A merely continuous compactly supported function need not have integrable
Fourier transform, so the hypothesis $F^{(m)}\in H^1_{\mathrm{loc}}$ is a real
additional condition. Accordingly, the general problem of replacing local
Wiener regularity by the minimal interior-$p$ hypotheses remains open here.

A bounded novelty search on 11 August 2026 queried the exact title and arXiv
identifier, the exact cusp symbol, ``remove this gap'', ``quantum Besov'', and
Wiener-space phrases. It found the current source but no later arXiv correction
or resolution of this observation. Novelty confidence is therefore moderate,
not definitive.

\paragraph{Human-review recommendation.}
Check the cutoff convention for $\in_{\mathrm{loc}}$ and the source's Fourier
normalization. Neither affects the argument. If confirmed, Remark~1.5's
specific example should be corrected, while its broader regularity question
should be retained.

\begin{thebibliography}{9}
\bibitem{ChenHong}
D. Chen and G. Hong,
\emph{The nonlinear estimates on quantum Besov spaces},
arXiv:2601.11934v2, 2026.
\end{thebibliography}

\end{document}
